\documentclass{article}

\usepackage[UTF8]{ctex}
\usepackage{amsmath, amssymb, amsthm, amsfonts}
\usepackage{geometry}
\usepackage{array, graphicx, pgfplots, tikz}
\usepackage{enumerate}
\usepackage{ulem}
\usepackage{float}

\newtheorem{theorem}{Theorem}[section]
\newtheorem{lemma}{Lemma}[section]
\newtheorem{definition}{Definition}[section]
\newtheorem{example}{Example}[section]

\numberwithin{equation}{section}

\pgfplotsset{compat=1.18}

\geometry{
    paperwidth=21cm,
    paperheight=29.7cm,
    top=2.5cm,
    bottom=2.5cm,
    left=2.5cm,
    right=2.5cm
}

\title{开题报告}
\author{陈昱景、张云航、曾沛然、李子昂}
\date{\today}

\begin{document}

\maketitle

\section{Excitation}

唬牌是民间很流行的一种扑克牌对战游戏，通常由2～3名玩家参与，使用一副标准扑克牌，其中大小王可被视为任意牌用于增加规则灵活性。游戏中，玩家需要依次倒扣出牌，并口头宣称“当前出牌的数量和牌面”，这些牌既可能是真实存在的，也可能是故意虚构的。对于下一位玩家而言，他需要在有限信息下做出判断：是继续相信上家并跟牌，还是质疑其真伪并发起揭牌，亦或者当局放弃出牌直接跳过。若最终某名玩家最先将手中全部牌出完，则该玩家获胜。唬牌的核心并不在于牌型大小，而在于对信息的解读与心理博弈：玩家需要不断推测对方是否在 bluff，从而决定是否跟随、揭穿或规避风险。

在现实生活中，唬牌通常需要亲自面对面交流，依赖玩家的语言表达、情绪变化与经验判断，因此更具互动性和娱乐性。然而，现阶段很多人只能在线下和朋友一起进行对战，缺少专门针对该类游戏的计算机对战系统，尤其是具备人机博弈功能的程序较少。为了让唬牌不仅能作为休闲娱乐，也能够成为计算机博弈与人工智能研究的案例，我们计划开发一套完整的唬牌游戏系统。该系统将采用 C++ 负责实现游戏核心逻辑，HTML + JavaScript 负责前端可视化交互，并嵌入 AI 出牌模块，实现人与机器人之间的对战。通过这套方案，我们能够将传统纸牌玩法数字化，并在规则引擎和智能决策方面展开系统化研究。

本课题的提出，正是基于唬牌游戏既具备强烈娱乐性，又具有明显信息不完全和心理对抗特征这一特点。其对弈过程不仅需要严格处理规则，还要求算法能够在有限信息条件下做出合理判断，因此对后端逻辑、前端界面和 AI 策略都提出了较高要求。我们希望通过本项目实现一个具备基本玩法、稳定流程和可扩展结构的唬牌系统，为后续进一步研究博弈类小游戏与 AI 决策提供技术基础和实验平台。

\section{Exploration}

在梳理唬牌的规则与开发需求后，我们将项目拆分为\textbf{后端}、\textbf{前端}和\textbf{人机AI} 三个核心模块。后端模块主要负责封装唬牌完整业务逻辑，包括牌堆初始化、玩家信息管理、卡牌分发、回合推进、动作合法性校验、对局状态更新和胜负判定等；前端模块负责实现界面展示与交互操作，采用 HTML 页面方案将玩家手牌、桌面牌、历史喊牌记录和当前状态直观呈现给用户；AI 模块则负责研究和设计机器人出牌策略，使其能够在有限信息下进行决策，并在对局中与用户进行真实对抗。

从系统结构上看，本课题本质上属于博弈类小游戏的系统设计与算法实现问题。一方面，需要保证游戏规则严谨、流程稳定和胜负判定准确，避免出现错误的出牌或判定；另一方面，还需要兼顾界面交互可用性和AI决策逻辑的合理性。例如，玩家在每一轮都需要判断上一轮说法是否可信，并结合历史记录、对手风格和手牌状态做出“跟牌”“揭牌”“过牌”的决策，这种决策过程具有典型的不完全信息博弈特征。同步设计这样一个系统，既能体现算法与工程的结合，也能够更好地研究人工智能在对战场景中的策略表现。

\begin{table}[H]
    \centering
    \renewcommand{\arraystretch}{1.2}
    \caption{分工情况}
    \label{tab:symbols}
    \begin{tabular}{ccc}
        \hline
        \textbf{姓名} & \textbf{分工} & \textbf{任务}\\
        \hline
        陈昱景 & 前端     & 设计游戏界面, 实现基本的人机交互\\
        李子昂 & 后端     & 设计游戏逻辑, 实现基本的游戏对局\\
        曾沛然 & 人机逻辑 & 设计人机逻辑, 实现基本的 AI 操作\\
        张云航 & 撰写报告 & 设计文档与报告, 制定任务方向和总结\\
        \hline
    \end{tabular}
\end{table}

此外，从项目实施角度看，第一阶段应以“规则引擎”与“对局控制器”为核心，优先完成完整的桌面逻辑实现，再逐步接入前端交互和机器人决策模块。该方法能够降低后期联调难度，提高系统构造的稳定性。最终形成的产品不仅具有娱乐属性，也能够作为一个较完整的博弈型小游戏框架，便于后续扩展多人联机、难度调节、战绩保存和复盘分析等功能。

\section{Enhancement}

本项目实施过程中，主要需要解决的关键问题包括：\textbf{后端规则的严谨性}、\textbf{前后端通信方式的稳定性}以及\textbf{AI决策策略的有效性}。其中，后端采用面向对象设计思路，用 C++ 中的类来分别表示牌、玩家、手牌、桌面状态和游戏控制器，能够实现模块化管理和状态隔离。这样可以保证每一轮操作都由一致的规则判断器处理，避免前端直接操控游戏状态带来的逻辑不一致问题。与此同时，为了支持浏览器端交互，前端拟采用 HTML + JavaScript 构建用户界面，并通过 WebSocket 或 HTTP 接口与后端建立通信，从而实现玩家动作的实时传输和游戏状态的同步刷新。

\begin{enumerate}
    \item \textbf{后端逻辑设计：} 需要实现牌堆生成、洗牌、发牌、出牌合法性校验、历史记录保存、轮次推进和胜负判定等核心逻辑。特别是在“喊牌是否真实”这一关键环节，系统必须严格判断当前牌型是否符合宣称数量和花色，且要正确处理大小王的特殊角色和边界情况，以确保出牌判定的一致性。

    \item \textbf{前端界面实现：} 前端需展示每位玩家的手牌数量、桌面已出牌、当前喊牌内容和历史记录，同时设计清晰的操作面板，让用户能够进行“跟牌”“揭牌”“过牌”等动作。界面应尽量简洁易用，并兼顾桌面端与浏览器展示效果，以便快速开展系统演示和测试。

    \item \textbf{人机 AI 设计：} AI 模块的关键目标是使机器人在有限信息下做出合理决策，而不是简单随机出牌。可通过历史出牌记录、相邻玩家的行为模式、当前手牌剩余情况和概率估计来模拟真实玩家的判断逻辑，并结合“跟牌-揭牌-过牌”三种决策方式进行权衡，以提高 AI 的挑战性与可玩性。

    \item \textbf{系统扩展与优化：} 后续还可以扩展为单人对战、多人在线、难度选择、游戏日志查询以及复盘分析等功能，从而提升系统的完整性和工程实用性。通过模块化设计，项目可以从最基础的原型逐步迭代成一个较成熟的博弈型游戏平台。
\end{enumerate}

综上所述，本项目的研究重点不仅是开发一款具备娱乐性的扑克牌游戏，更是在规则设计、前后端协同和智能决策等方面形成一套完整的工程化实现方案。通过唬牌这一典型的心理博弈游戏，我们能够检验软件架构设计、交互体验和 AI 对抗性在实际项目中的协同效果，从而为后续更复杂的对战型游戏开发奠定基础。

\end{document}
\end{document}
